We can see how efficiently complex computations can be performed using tensor networks. We now introduce other identities involving random Gaussian matrices.
\begin{enumerate}
\item  The first identity is a useful consequence of \emph{Isserlis'} theorem~\citep{isserlis1918formula}.
\begin{theorem}\label{thm:isserlis}
Let $\ab\in\RR^{n}$ be a random vector whose elements are drawn i.i.d. from the normal distribution with mean zero and variance one. Then
\begin{align*}
\EE(\ab^{\outprod 4}) 
&= 
\EE\left(\begin{tikzpicture}[baseline=\baseline]
   \tikzset{tensor/.style = {minimum size = 0.5cm,shape = circle,thick,draw=black,inner sep = 0pt}, edge/.style = {thick,line width=.4mm},every loop/.style={}}
    \def\y{0.25}
    \def\x{0}
    \node[tensor,fill=red!20!blue!20!white] (a1) at (\x,\y){\scalebox{0.85}{$\ab$}}; 
    \node[tensor,fill=red!20!blue!20!white] (a2) at (\x+1,\y){\scalebox{0.85}{$\ab$}}; 
    \node[tensor,fill=red!20!blue!20!white] (a3) at (\x+2,\y){\scalebox{0.85}{$\ab$}}; 
    \node[tensor,fill=red!20!blue!20!white] (a4) at (\x+3,\y){\scalebox{0.85}{$\ab$}}; 
    \draw[edge](a1) -- (\x,\y-0.85);
    \draw[edge](a2) -- (\x+1,\y-0.85);
    \draw[edge](a3) -- (\x+2,\y-0.85);
    \draw[edge](a4) -- (\x+3,\y-0.85);
\end{tikzpicture}
\right)\\
&=
\begin{tikzpicture}[baseline=\baseline]
    \tikzset{tensor/.style = {minimum size = 0.5cm,shape = circle,thick,draw=black,inner sep = 0pt}, edge/.style = {thick,line width=.4mm},every loop/.style={}}
    \def\y{0}
    \def\x{0}
    \draw[edge] (\x+0.15,\y-0.1) to [out=100,in=80] (\x+1.85,\y-0.1);
    \draw[edge] (\x+2,\y-0.1) to [out=100,in=80] (\x+3.65,\y-0.1);
\end{tikzpicture}
+
\begin{tikzpicture}[baseline=\baseline]
    \tikzset{tensor/.style = {minimum size = 0.5cm,shape = circle,thick,draw=black,inner sep = 0pt}, edge/.style = {thick,line width=.4mm},every loop/.style={}}
    \def\y{0}
    \def\x{0}
    \draw[edge] (\x+0.15,\y-0.1) to [out=100,in=80] (\x+1.85,\y-0.1);
    \draw[edge] (\x+1.45,\y-0.1) to [out=100,in=80] (\x+3,\y-0.1);
\end{tikzpicture}
+
    \begin{tikzpicture}[baseline=\baseline]
    \tikzset{tensor/.style = {minimum size = 0.5cm,shape = circle,thick,draw=black,inner sep = 0pt}, edge/.style = {thick,line width=.4mm},every loop/.style={}}
    \def\y{0}
    \def\x{0}
    \draw[edge] (\x,\y-0.1) to [out=100,in=80] (\x+2,\y-0.1);
    \draw[edge] (\x+0.15,\y-0.1) to [out=100,in=80] (\x+1.85,\y-0.1);
\end{tikzpicture}
\end{align*}
where $\ab^{\outprod 4}$ denotes the outer product of $\ab$ with itself 4 times. In the last equality, each term in the sum corresponds to a different reshaping of the identity matrix.

This result can be extended to an arbitrary covariance structure: if $\ab$ is drawn from the multivariate normal $\mathcal N(\mathbf 0, \Sigmab )$, then
\begin{align*}
\EE(\ab^{\outprod 4})=
\begin{tikzpicture}[baseline=\baseline]
    \tikzset{tensor/.style = {minimum size = 0.4cm,shape = circle,thick,draw=black,inner sep = 0pt}, edge/.style = {thick,line width=.4mm},every loop/.style={}}
    \def\y{0}`
    \def\x{0}
    \draw[edge] (\x+0.15,\y-0.1) to [out=100,in=80] (\x+1.85,\y-0.1);
    \draw[edge] (\x+2,\y-0.1) to [out=100,in=80] (\x+3.65,\y-0.1);
    \node[tensor,fill=red!50!white] (A) at (\x+0.9,\y+0.35){\scalebox{0.85}{$\Sigmab$}};
    \node[tensor,fill=red!50!white] (A) at (\x+2.8,\y+0.35){\scalebox{0.85}{$\Sigmab$}};
\end{tikzpicture}
+
\begin{tikzpicture}[baseline=\baseline]
    \tikzset{tensor/.style = {minimum size = 0.4cm,shape = circle,thick,draw=black,inner sep = 0pt}, edge/.style = {thick,line width=.4mm},every loop/.style={}}
    \def\y{0}`
    \def\x{0}
    \draw[edge] (\x+0.15,\y-0.1) to [out=100,in=80] (\x+1.85,\y-0.1);
    \draw[edge] (\x+1.45,\y-0.1) to [out=100,in=80] (\x+3,\y-0.1);
    \node[tensor,fill=red!50!white] (A) at (\x+0.9,\y+0.35){\scalebox{0.85}{$\Sigmab$}};
    \node[tensor,fill=red!50!white] (A) at (\x+2.25,\y+0.35){\scalebox{0.85}{$\Sigmab$}};
\end{tikzpicture}
+
\begin{tikzpicture}[baseline=\baseline]
    \tikzset{tensor/.style = {minimum size = 0.4cm,shape = circle,thick,draw=black,inner sep = 0pt}, edge/.style = {thick,line width=.4mm},every loop/.style={}}
    \def\y{0.3}`
    \def\x{0}
    \draw[edge] (\x,\y-0.1) to [out=100,in=80] (\x+2,\y-0.1);
    \draw[edge] (\x+0.15,\y-0.45) to [out=100,in=80] (\x+1.85,\y-0.45);
    \node[tensor,fill=red!50!white] (A) at (\x+0.9,\y+0.45){\scalebox{0.85}{$\Sigmab$}};
    \node[tensor,fill=red!50!white] (A) at (\x+0.95,\y){\scalebox{0.85}{$\Sigmab$}};
\end{tikzpicture} .
\end{align*}


\begin{proof}
    We show the result for the case where the covariance is the identity. The proof extends straightforwardly to the case of an arbitrary covariance matrix. 
    
    Element-wise, by using Isserlis’
    theorem and using the fact that $\ab\sim\mathcal{N}(\mathbf{0},\Ib)$, for $i_1,\cdots,i_4\in[n]$, we have
\begin{align*}
(\EE(&\ab^{\outprod 4}))_{i_1,i_2,i_3,i_4} 
=
    \EE(\ab_{i_1}\ab_{i_2}\ab_{i_3}\ab_{i_4})\\
&= 
    \EE(\ab_{i_1}\ab_{i_2})\EE(\ab_{i_3}\ab_{i_4}) + \EE(\ab_{i_1}\ab_{i_3})\EE(\ab_{i_2}\ab_{i_4})
    +\EE(\ab_{i_1}\ab_{i_4})\EE(\ab_{i_2}\ab_{i_3})\\
&=
    \delta_{i_1i_2}\delta_{i_3i_4} + \delta_{i_1i_3}\delta_{i_2i_4} + \delta_{i_1i_4}\delta_{i_2i_3},
\end{align*}
which in tensor network notation gives:
\begin{align*}
\EE(\ab^{\outprod 4})_{i_1i_2i_3i_4}
=
\begin{tikzpicture}[baseline=\baseline]
    \tikzset{tensor/.style = {minimum size = 0.5cm,shape = circle,thick,draw=black,inner sep = 0pt}, edge/.style = {thick,line width=.4mm},every loop/.style={}}
    \def\y{0}
    \def\x{0}
    \draw[edge] (\x+0.15,\y-0.1) to [out=100,in=80] (\x+1.85,\y-0.1);
    \draw[edge] (\x+2,\y-0.1) to [out=100,in=80] (\x+3.65,\y-0.1);
    \node[draw=none] () at (\x+0.15,\y-0.25) {\scalebox{0.7}{\indcolor{$i_1$}}};
    \node[draw=none] () at (\x+1.85,\y-0.25) {\scalebox{0.7}{\indcolor{$i_2$}}};
    \node[draw=none] () at (\x+2.15,\y-0.25) {\scalebox{0.7}{\indcolor{$i_3$}}};
    \node[draw=none] () at (\x+3.65,\y-0.25) {\scalebox{0.7}{\indcolor{$i_4$}}};
\end{tikzpicture}
+
\begin{tikzpicture}[baseline=\baseline]
    \tikzset{tensor/.style = {minimum size = 0.5cm,shape = circle,thick,draw=black,inner sep = 0pt}, edge/.style = {thick,line width=.4mm},every loop/.style={}}
    \def\y{0}
    \def\x{0}
    \draw[edge] (\x+0.15,\y-0.1) to [out=100,in=80] (\x+1.85,\y-0.1);
    \draw[edge] (\x+1.45,\y-0.1) to [out=100,in=80] (\x+3,\y-0.1);
    \node[draw=none] () at (\x+0.15,\y-0.25) {\scalebox{0.7}{\indcolor{$i_1$}}};
    \node[draw=none] () at (\x+3,\y-0.25) {\scalebox{0.7}{\indcolor{$i_4$}}};
    \node[draw=none] () at (\x+1.85,\y-0.25) {\scalebox{0.7}{\indcolor{$i_3$}}};
    \node[draw=none] () at (\x+1.45,\y-0.25) {\scalebox{0.7}{\indcolor{$i_2$}}};
\end{tikzpicture}
+
    \begin{tikzpicture}[baseline=\baseline]
    \tikzset{tensor/.style = {minimum size = 0.5cm,shape = circle,thick,draw=black,inner sep = 0pt}, edge/.style = {thick,line width=.4mm},every loop/.style={}}
    \def\y{0}
    \def\x{0}
    \draw[edge] (\x,\y-0.1) to [out=100,in=80] (\x+2,\y-0.1);
    \draw[edge] (\x+0.15,\y-0.1) to [out=100,in=80] (\x+1.85,\y-0.1);
    \node[draw=none] () at (\x,\y-0.25) {\scalebox{0.7}{\indcolor{$i_1$}}};
    \node[draw=none] () at (\x+2.15,\y-0.25) {\scalebox{0.7}{\indcolor{$i_4$}}};
    \node[draw=none] () at (\x+0.35,\y-0.25) {\scalebox{0.7}{\indcolor{$i_2$}}};
    \node[draw=none] () at (\x+1.7,\y-0.25) {\scalebox{0.7}{\indcolor{$i_3$}}};
\end{tikzpicture}.
\end{align*}
\end{proof}
\end{theorem}
We can also extend the result of \emph{Isserlis' theorem} to the case where the random variable takes the form of a higher-order tensor. For example, for a third-order tensor $\At\in\RR^{d_1\times d_2\times d_3}$, we can write
\begin{align*}
    (\EE(\At^{\outprod 4}))_{i_1j_1r_1i_2j_2r_2i_3j_3r_3i_4j_4r_4}
&=\\
&\hspace{-3cm}\begin{tikzpicture}[baseline=\baseline]
    \tikzset{tensor/.style = {minimum size = 0.5cm,shape = circle,thick,draw=black,inner sep = 0pt}, edge/.style = {thick,line width=.4mm},every loop/.style={}}
    \def\x{0}
    \def\y{0}
    \draw[edge] (\x+0.15,\y-0.1) to [out=100,in=80] (\x+1.85,\y-0.1);
    \draw[edge] (\x+2,\y-0.1) to [out=100,in=80] (\x+3.65,\y-0.1);
    \node[draw=none] () at (\x+0.15,\y-0.25) {\scalebox{0.7}{\indcolor{$j_1$}}};
    \node[draw=none] () at (\x+1.85,\y-0.25) {\scalebox{0.7}{\indcolor{$j_2$}}};
    \node[draw=none] () at (\x+2.15,\y-0.25) {\scalebox{0.7}{\indcolor{$j_3$}}};
    \node[draw=none] () at (\x+3.65,\y-0.25) {\scalebox{0.7}{\indcolor{$j_4$}}};
    \draw[edge] (\x+0.15,\y+1) to [out=100,in=80] (\x+1.85,\y+1);
    \draw[edge] (\x+2,\y+1) to [out=100,in=80] (\x+3.65,\y+1);
    \node[draw=none] () at (\x+0.15,\y+0.9) {\scalebox{0.7}{\indcolor{$i_1$}}};
    \node[draw=none] () at (\x+1.85,\y+0.9) {\scalebox{0.7}{\indcolor{$i_2$}}};
    \node[draw=none] () at (\x+2.15,\y+0.9) {\scalebox{0.7}{\indcolor{$i_3$}}};
    \node[draw=none] () at (\x+3.65,\y+0.9) {\scalebox{0.7}{\indcolor{$i_4$}}};
    \draw[edge] (\x+0.15,\y-1) to [out=100,in=80] (\x+1.85,\y-1);
    \draw[edge] (\x+2,\y-1) to [out=100,in=80] (\x+3.65,\y-1);
    \node[draw=none] () at (\x+0.15,\y-1.15) {\scalebox{0.7}{\indcolor{$r_1$}}};
    \node[draw=none] () at (\x+1.85,\y-1.15) {\scalebox{0.7}{\indcolor{$r_2$}}};
    \node[draw=none] () at (\x+2.15,\y-1.15) {\scalebox{0.7}{\indcolor{$r_3$}}};
    \node[draw=none] () at (\x+3.65,\y-1.15) {\scalebox{0.7}{\indcolor{$r_4$}}};
\end{tikzpicture}
+
\begin{tikzpicture}[baseline=\baseline]
    \tikzset{tensor/.style = {minimum size = 0.5cm,shape = circle,thick,draw=black,inner sep = 0pt}, edge/.style = {thick,line width=.4mm},every loop/.style={}}
    \def\x{0}
    \def\y{0}
    \draw[edge] (\x+0.15,\y-0.1) to [out=100,in=80] (\x+1.85,\y-0.1);
    \draw[edge] (\x+1.45,\y-0.1) to [out=100,in=80] (\x+3,\y-0.1);
    \node[draw=none] () at (\x+0.15,\y-0.25) {\scalebox{0.7}{\indcolor{$j_1$}}};
    \node[draw=none] () at (\x+3,\y-0.25) {\scalebox{0.7}{\indcolor{$j_4$}}};
    \node[draw=none] () at (\x+1.85,\y-0.25) {\scalebox{0.7}{\indcolor{$j_3$}}};
    \node[draw=none] () at (\x+1.45,\y-0.25) {\scalebox{0.7}{\indcolor{$j_2$}}};
    \draw[edge] (\x+0.15,\y+1) to [out=100,in=80] (\x+1.85,\y+1);
    \draw[edge] (\x+1.45,\y+1) to [out=100,in=80] (\x+3,\y+1);
    \node[draw=none] () at (\x+0.15,\y+0.9) {\scalebox{0.7}{\indcolor{$i_1$}}};
    \node[draw=none] () at (\x+3,\y+0.9) {\scalebox{0.7}{\indcolor{$i_4$}}};
    \node[draw=none] () at (\x+1.85,\y+0.9) {\scalebox{0.7}{\indcolor{$i_3$}}};
    \node[draw=none] () at (\x+1.45,\y+0.9) {\scalebox{0.7}{\indcolor{$i_2$}}};
    \draw[edge] (\x+0.15,\y-1) to [out=100,in=80] (\x+1.85,\y-1);
    \draw[edge] (\x+1.45,\y-1) to [out=100,in=80] (\x+3,\y-1);
    \node[draw=none] () at (\x+0.15,\y-1.15) {\scalebox{0.7}{\indcolor{$r_1$}}};
    \node[draw=none] () at (\x+3,\y-1.15) {\scalebox{0.7}{\indcolor{$r_4$}}};
    \node[draw=none] () at (\x+1.85,\y-1.15) {\scalebox{0.7}{\indcolor{$r_3$}}};
    \node[draw=none] () at (\x+1.45,\y-1.15) {\scalebox{0.7}{\indcolor{$r_2$}}};
\end{tikzpicture}
+
    \begin{tikzpicture}[baseline=\baseline]
    \tikzset{tensor/.style = {minimum size = 0.5cm,shape = circle,thick,draw=black,inner sep = 0pt}, edge/.style = {thick,line width=.4mm},every loop/.style={}}
    \def\x{0}
    \def\y{0}
    \draw[edge] (\x,\y+1) to [out=100,in=80] (\x+2,\y+1);
    \draw[edge] (\x+0.15,\y+1) to [out=100,in=80] (\x+1.85,\y+1);
    \node[draw=none] () at (\x,\y+0.9) {\scalebox{0.7}{\indcolor{$i_1$}}};
    \node[draw=none] () at (\x+2.15,\y+0.9) {\scalebox{0.7}{\indcolor{$i_4$}}};
    \node[draw=none] () at (\x+0.35,\y+0.9) {\scalebox{0.7}{\indcolor{$i_2$}}};
    \node[draw=none] () at (\x+1.7,\y+0.9) {\scalebox{0.7}{\indcolor{$i_3$}}};
    \draw[edge] (\x,\y-0.1) to [out=100,in=80] (\x+2,\y-0.1);
    \draw[edge] (\x+0.15,\y-0.1) to [out=100,in=80] (\x+1.85,\y-0.1);
    \node[draw=none] () at (\x,\y-0.25) {\scalebox{0.7}{\indcolor{$j_1$}}};
    \node[draw=none] () at (\x+2.15,\y-0.25) {\scalebox{0.7}{\indcolor{$j_4$}}};
    \node[draw=none] () at (\x+0.35,\y-0.25) {\scalebox{0.7}{\indcolor{$j_2$}}};
    \node[draw=none] () at (\x+1.7,\y-0.25) {\scalebox{0.7}{\indcolor{$j_3$}}};
    \draw[edge] (\x,\y-1) to [out=100,in=80] (\x+2,\y-1);
    \draw[edge] (\x+0.15,\y-1) to [out=100,in=80] (\x+1.85,\y-1);
    \node[draw=none] () at (\x,\y-1.25) {\scalebox{0.7}{\indcolor{$r_1$}}};
    \node[draw=none] () at (\x+2.15,\y-1.25) {\scalebox{0.7}{\indcolor{$r_4$}}};
    \node[draw=none] () at (\x+0.35,\y-1.25) {\scalebox{0.7}{\indcolor{$r_2$}}};
    \node[draw=none] () at (\x+1.7,\y-1.25) {\scalebox{0.7}{\indcolor{$r_3$}}};
\end{tikzpicture},
\end{align*}
for $i_1,\cdots,i_4\in[d_1], j_1,\cdots,j_4\in [d_2]$ and $r_1,\cdots,r_4\in[d_3]$.
\item Using the above results, we can easily derive the expectation of  $(\Ab^\ts\Ab)^{\outprod 2}$ when $\Ab$ is a random matrix:
\begin{proposition}\label{prop:kronecker-exp}
    Let $\Ab\in\RR^{m\times n}$ be a random matrix whose elements are drawn i.i.d from the standard normal distribution. Then
\begin{align*}
\EE((\Ab^\ts\Ab)^{\outprod 2})) 
=
m^2\begin{tikzpicture}[baseline=\baseline]
    \tikzset{tensor/.style = {minimum size = 0.5cm,shape = circle,thick,draw=black,inner sep = 0pt}, edge/.style = {thick,line width=.4mm},every loop/.style={}}
    \def\y{0}
    \def\x{0}
    \draw[edge] (\x+0.15,\y-0.1) to [out=100,in=80] (\x+1.85,\y-0.1);
    \draw[edge] (\x+2,\y-0.1) to [out=100,in=80] (\x+3.65,\y-0.1);
    \node[draw=none] () at (\x+1,\y+0.5) {\scalebox{0.7}{\dimcolor{$n$}}};
    \node[draw=none] () at (\x+2.85,\y+0.5) {\scalebox{0.7}{\dimcolor{$n$}}};
\end{tikzpicture}
+m\begin{tikzpicture}[baseline=\baseline]
    \tikzset{tensor/.style = {minimum size = 0.5cm,shape = circle,thick,draw=black,inner sep = 0pt}, edge/.style = {thick,line width=.4mm},every loop/.style={}}
    \def\y{0}
    \def\x{0}
    \draw[edge] (\x+0.15,\y-0.1) to [out=100,in=80] (\x+1.85,\y-0.1);
    \draw[edge] (\x+1.45,\y-0.1) to [out=100,in=80] (\x+3,\y-0.1);
    \node[draw=none] () at (\x+1,\y+0.5) {\scalebox{0.7}{\dimcolor{$n$}}};
    \node[draw=none] () at (\x+2.25,\y+0.5) {\scalebox{0.7}{\dimcolor{$n$}}};
\end{tikzpicture}
+
m\begin{tikzpicture}[baseline=\baseline]
    \tikzset{tensor/.style = {minimum size = 0.5cm,shape = circle,thick,draw=black,inner sep = 0pt}, edge/.style = {thick,line width=.4mm},every loop/.style={}}
    \def\y{0}
    \def\x{0}
    \draw[edge] (\x,\y-0.1) to [out=100,in=80] (\x+2,\y-0.1);
    \draw[edge] (\x+0.15,\y-0.1) to [out=100,in=80] (\x+1.85,\y-0.1);
    \node[draw=none] () at (\x+1,\y+0.65) {\scalebox{0.7}{\dimcolor{$n$}}};
    \node[draw=none] () at (\x+1,\y+0.25) {\scalebox{0.7}{\dimcolor{$n$}}};
\end{tikzpicture}.
\end{align*}
\begin{proof}
By applying Theorem~\ref{thm:isserlis} to matrices and using Proposition~\ref{prop:wishart-mean}, we can derive the desired expression. First, from Theorem~\ref{thm:isserlis} we have
\begin{align*}
\EE(\Ab^{\outprod4})_{i_1j_1i_2j_2i_3j_3i_4j_4}
&= 
\EE\left(\begin{tikzpicture}[baseline=\baseline]
    \tikzset{tensor/.style = {minimum size = 0.5cm,shape = circle,thick,draw=black,inner sep = 0pt}, edge/.style = {thick,line width=.4mm},every loop/.style={}}
    \def\y{0.25}
    \def\x{0}
    \node[tensor,fill=green!50!blue!20!white] (A1) at (\x,\y){\scalebox{0.85}{$\Ab$}};
    \node[tensor,fill=green!50!blue!20!white] (A2) at (\x+1,\y){\scalebox{0.85}{$\Ab$}};
    \node[tensor,fill=green!50!blue!20!white] (A3) at (\x+2,\y){\scalebox{0.85}{$\Ab$}};
    \node[tensor,fill=green!50!blue!20!white] (A4) at (\x+3,\y){\scalebox{0.85}{$\Ab$}};
    \draw[edge](\x-0.15,\y-0.18) -- (\x-0.15,\y-0.65)node [below]{\scalebox{0.7}{\indcolor{$i_1$}}};
    \draw[edge](\x+0.15,\y-0.18) -- (\x+0.15,\y-0.65)node [below]{\scalebox{0.7}{\indcolor{$j_1$}}};
    \draw[edge](\x+0.85,\y-0.18) -- (\x+0.85,\y-0.65)node [below]{\scalebox{0.7}{\indcolor{$i_2$}}};
    \draw[edge](\x+1.15,\y-0.18) -- (\x+1.15,\y-0.65)node [below]{\scalebox{0.7}{\indcolor{$j_2$}}};
    \draw[edge](\x+1.85,\y-0.18) -- (\x+1.85,\y-0.65)node [below]{\scalebox{0.7}{\indcolor{$i_3$}}};
    \draw[edge](\x+2.15,\y-0.18) -- (\x+2.15,\y-0.65)node [below]{\scalebox{0.7}{\indcolor{$j_3$}}};
    \draw[edge](\x+2.85,\y-0.18) -- (\x+2.85,\y-0.65)node [below]{\scalebox{0.7}{\indcolor{$i_4$}}};
    \draw[edge](\x+3.15,\y-0.18) -- (\x+3.15,\y-0.65)node [below]{\scalebox{0.7}{\indcolor{$j_4$}}};
     \node[draw=none] () at (\x-0.35,\y-0.4) {\scalebox{0.7}{\dimcolor{$n$}}};
    \node[draw=none] () at (\x+0.35,\y-0.4) {\scalebox{0.7}{\dimcolor{$m$}}};
    \node[draw=none] () at (\x+0.65,\y-0.4) {\scalebox{0.7}{\dimcolor{$n$}}};
    \node[draw=none] () at (\x+1.35,\y-0.4) {\scalebox{0.7}{\dimcolor{$m$}}};
    \node[draw=none] () at (\x+1.65,\y-0.4) {\scalebox{0.7}{\dimcolor{$n$}}};
    \node[draw=none] () at (\x+2.35,\y-0.4) {\scalebox{0.7}{\dimcolor{$m$}}};
    \node[draw=none] () at (\x+2.65,\y-0.4) {\scalebox{0.7}{\dimcolor{$n$}}};
    \node[draw=none] () at (\x+3.35,\y-0.4) {\scalebox{0.7}{\dimcolor{$m$}}};
\end{tikzpicture}\right)\\
&\hspace{-1.5cm}=  
\begin{tikzpicture}[baseline=\baseline]
    \tikzset{tensor/.style = {minimum size = 0.5cm,shape = circle,thick,draw=black,inner sep = 0pt}, edge/.style = {thick,line width=.4mm},every loop/.style={}}
    \def\y{0.5}
    \def\x{0}
    \draw[edge] (\x+0.15,\y-0.1) to [out=100,in=80] (\x+1.85,\y-0.1);
    \draw[edge] (\x+2,\y-0.1) to [out=100,in=80] (\x+3.65,\y-0.1);
    \node[draw=none] () at (\x+0.15,\y-0.25) {\scalebox{0.7}{\indcolor{$i_1$}}};
    \node[draw=none] () at (\x+1.85,\y-0.25) {\scalebox{0.7}{\indcolor{$i_2$}}};
    \node[draw=none] () at (\x+2.15,\y-0.25) {\scalebox{0.7}{\indcolor{$i_3$}}};
    \node[draw=none] () at (\x+3.85,\y-0.25) {\scalebox{0.7}{\indcolor{$i_4$}}};
    \draw[edge](\x+0.15,\y-0.85) to [out=100,in=80] (\x+1.85,\y-0.85);
    \draw[edge] (\x+2,\y-0.85) to [out=100,in=80] (\x+3.65,\y-0.85);
    \node[draw=none] () at (\x+0.1,\y-0.95) {\scalebox{0.7}{\indcolor{$j_1$}}};
    \node[draw=none] () at (\x+1.8, \y-0.95) {\scalebox{0.7}{\indcolor{$j_2$}}};
    \node[draw=none] () at (\x+2.15,\y-0.95) {\scalebox{0.7}{\indcolor{$j_3$}}};
    \node[draw=none] () at (\x+3.85,\y-0.95) {\scalebox{0.7}{\indcolor{$j_4$}}};
\end{tikzpicture}
+
\begin{tikzpicture}[baseline=\baseline]
    \tikzset{tensor/.style = {minimum size = 0.5cm,shape = circle,thick,draw=black,inner sep = 0pt}, edge/.style = {thick,line width=.4mm},every loop/.style={}}
    \def\y{0.5}
    \def\x{0}
    \draw[edge] (\x+0.15,\y-0.1) to [out=100,in=80] (\x+1.85,\y-0.1);
    \draw[edge] (\x+1.45,\y-0.1) to [out=100,in=80] (\x+3,\y-0.1);
    \node[draw=none] () at (\x+0.15,\y-0.25) {\scalebox{0.7}{\indcolor{$i_1$}}};
    \node[draw=none] () at (\x+3,\y-0.25) {\scalebox{0.7}{\indcolor{$i_4$}}};
    \node[draw=none] () at (\x+1.85,\y-0.25) {\scalebox{0.7}{\indcolor{$i_3$}}};
    \node[draw=none] () at (\x+1.45,\y-0.25) {\scalebox{0.7}{\indcolor{$i_2$}}};
    \draw[edge] (\x+0.15,\y-0.85) to [out=100,in=80] (\x+1.85,\y-0.85);
    \draw[edge] (\x+1.45,\y-0.85) to [out=100,in=80] (\x+3,\y-0.85);
    \node[draw=none] () at (\x+0.1,\y-0.95) {\scalebox{0.7}{\indcolor{$j_1$}}};
    \node[draw=none] () at (\x+2.95,\y-0.95) {\scalebox{0.7}{\indcolor{$j_4$}}};
    \node[draw=none] () at (\x+1.8,\y-0.95) {\scalebox{0.7}{\indcolor{$j_3$}}};
    \node[draw=none] () at (\x+1.4,\y-0.95) {\scalebox{0.7}{\indcolor{$j_2$}}};
\end{tikzpicture}
+
 \begin{tikzpicture}[baseline=\baseline]
    \tikzset{tensor/.style = {minimum size = 0.5cm,shape = circle,thick,draw=black,inner sep = 0pt}, edge/.style = {thick,line width=.4mm},every loop/.style={}}
    \def\y{0.5}
    \def\x{0}
    \draw[edge] (\x,\y-0.1) to [out=100,in=80] (\x+2,\y-0.1);
    \draw[edge] (\x+0.15,\y-0.1) to [out=100,in=80] (\x+1.85,\y-0.1);
    \node[draw=none] () at (\x,\y-0.25) {\scalebox{0.7}{\indcolor{$i_1$}}};
    \node[draw=none] () at (\x+2.15,\y-0.25) {\scalebox{0.7}{\indcolor{$i_4$}}};
    \node[draw=none] () at (\x+0.35,\y-0.23) {\scalebox{0.7}{\indcolor{$i_2$}}};
    \node[draw=none] () at (\x+1.7,\y-0.25) {\scalebox{0.7}{\indcolor{$i_3$}}};
    \draw[edge] (\x,\y-0.85) to [out=100,in=80] (\x+2,\y-0.85);
    \draw[edge] (\x+0.15,\y-0.85) to [out=100,in=80] (\x+1.85,\y-0.85);
    \node[draw=none] () at (\x-0.03,\y-0.95) {\scalebox{0.7}{\indcolor{$j_1$}}};
    \node[draw=none] () at (\x+2.15,\y-0.95) {\scalebox{0.7}{\indcolor{$j_4$}}};
    \node[draw=none] () at (\x+0.35,\y-0.95) {\scalebox{0.7}{\indcolor{$j_2$}}};
    \node[draw=none] () at (\x+1.7,\y-0.95) {\scalebox{0.7}{\indcolor{$j_3$}}};
\end{tikzpicture},
\end{align*}
